% ============================================
% СПИСОК СОКРАЩЕНИЙ И УСЛОВНЫХ ОБОЗНАЧЕНИЙ
% Оформлено по ГОСТ 7.32 п.4.8: столбцом без знаков препинания в конце
% строк; слева без абзацного отступа в алфавитном порядке — сокращения,
% справа через тире — их расшифровка. Латинские сокращения предшествуют
% кириллическим.
% ============================================

\begingroup
\setlength{\parindent}{0pt}
\setlength{\parskip}{0pt}
\newcommand{\abbr}[2]{%
	\noindent #1~-- #2\par
	\vspace{0.3em}
}

\abbr{AAD}{Acoustic Activity Detection, обнаружение акустической активности}
\abbr{ASIC}{Application-Specific Integrated Circuit, специализированная интегральная схема}
\abbr{BLE}{Bluetooth Low Energy, беспроводной интерфейс малого энергопотребления}
\abbr{CMSIS-NN}{Cortex Microcontroller Software Interface Standard for Neural Networks, библиотека нейросетевых ядер Arm}
\abbr{CNN}{Convolutional Neural Network, свёрточная нейронная сеть}
\abbr{CRNN}{Convolutional Recurrent Neural Network, свёрточно-рекуррентная нейронная сеть}
\abbr{DMA}{Direct Memory Access, прямой доступ к памяти}
\abbr{DS-CNN}{Depthwise Separable Convolutional Neural Network, свёрточная нейронная сеть с разделимыми по глубине свёртками}
\abbr{DSP}{Digital Signal Processor, цифровой сигнальный процессор}
\abbr{ESP-IDF}{Espressif IoT Development Framework, фреймворк разработки для устройств Espressif}
\abbr{ESP-NN}{Espressif Neural Network kernels, библиотека ядер нейросетевого инференса}
\abbr{ESP-SR}{Espressif Speech Recognition, фреймворк распознавания речи Espressif}
\abbr{EXT0, EXT1}{External wakeup sources, аппаратные источники пробуждения микроконтроллера из режима глубокого сна}
\abbr{FAR}{False Acceptance Rate, доля ложных срабатываний}
\abbr{FN}{False Negative, ложноотрицательное срабатывание}
\abbr{FP32}{32-bit floating-point, представление числа с плавающей точкой одинарной точности}
\abbr{FPGA}{Field-Programmable Gate Array, программируемая логическая интегральная схема}
\abbr{FRR}{False Rejection Rate, доля ложных отказов}
\abbr{GPIO}{General-Purpose Input/Output, универсальный порт ввода-вывода}
\abbr{I2C}{Inter-Integrated Circuit, двухпроводная последовательная шина}
\abbr{I2S}{Inter-IC Sound, последовательный цифровой аудиоинтерфейс}
\abbr{INT8}{8-bit signed integer, 8-битное целое со знаком}
\abbr{IoT}{Internet of Things, интернет вещей}
\abbr{KWS}{Keyword Spotting, распознавание ключевых слов}
\abbr{LDO}{Low-Dropout regulator, линейный стабилизатор напряжения с малым падением}
\abbr{MAC}{Multiply-Accumulate, операция умножения с накоплением}
\abbr{MCU}{Microcontroller Unit, микроконтроллер}
\abbr{MEMS}{Microelectromechanical Systems, микроэлектромеханические системы}
\abbr{MFCC}{Mel-Frequency Cepstral Coefficients, мел-частотные кепстральные коэффициенты}
\abbr{MLPerf}{Machine Learning Performance benchmark, эталонный набор тестов производительности машинного обучения}
\abbr{PSRAM}{Pseudo-Static Random Access Memory, псевдостатическая оперативная память}
\abbr{PTQ}{Post-Training Quantization, посттренировочная квантизация}
\abbr{QAT}{Quantization-Aware Training, квантизация с учётом при обучении}
\abbr{ReLU}{Rectified Linear Unit, выпрямленная линейная функция активации}
\abbr{RNN}{Recurrent Neural Network, рекуррентная нейронная сеть}
\abbr{RTC}{Real-Time Clock, часы реального времени}
\abbr{SIMD}{Single Instruction Multiple Data, одна инструкция над множеством данных}
\abbr{SRAM}{Static Random Access Memory, статическая оперативная память}
\abbr{STE}{Straight-Through Estimator, метод сквозного прохождения градиента}
\abbr{TC-ResNet}{Temporal Convolution Residual Network, остаточная сеть с временными свёртками}
\abbr{TFLM}{TensorFlow Lite Micro, среда исполнения нейронных сетей для микроконтроллеров}
\abbr{UART}{Universal Asynchronous Receiver-Transmitter, универсальный асинхронный приёмопередатчик}
\abbr{ULP}{Ultra-Low Power coprocessor, сопроцессор сверхмалой мощности}
\abbr{USB}{Universal Serial Bus, универсальная последовательная шина}
\abbr{VAD}{Voice Activity Detector, детектор голосовой активности}
\abbr{Xtensa LX7}{архитектура ядра процессоров ESP32-S3}
\abbr{БПФ}{быстрое преобразование Фурье}
\abbr{ДКП}{дискретное косинусное преобразование}
\abbr{ДПФ}{дискретное преобразование Фурье}
\abbr{п.п.}{процентные пункты}
\abbr{ПЛИС}{программируемая логическая интегральная схема}
\abbr{ЦПУ}{центральное процессорное устройство}
\endgroup
